% paper/sections/01b_classification_roadmap.tex
% ═══════════════════════════════════════════════════════════════════════════════
\subsection{界面捕捉法の分類と本稿の手法の位置付け}
% ═══════════════════════════════════════════════════════════════════════════════

気液界面を数値的に追跡する手法は，大きく以下の3種類に分類される．
それぞれに保存性・幾何精度・界面鮮鋭性のトレードオフがある．
本稿の Conservative Level Set（CLS）法は，保存形移流を主変数 $\psi$ で行い，
幾何量を補助距離場 $\phi$ で評価することで，このトレードオフを明示的に管理する立場をとる．

\begin{description}
  \item[\textbf{VOF（Volume of Fluid）法}~\cite{HirtNichols1981}]
        各セル内の液相体積分率 $F \in [0,1]$ を輸送変数とする．
        保存形の輸送方程式を解くため質量保存性に優れるが，
        界面の幾何学的情報（法線 $\hat{\bm n}_{lg}$，曲率 $\kappa_{lg}$）を $F$ から
        再構成する精度が限られる．寄生流れが生じやすい．

  \item[\textbf{標準 Level Set（LS）法}~\cite{OsherSethian1988}]
        界面からの符号付き距離関数 $\phi$（$|\bnabla\phi|=1$）を輸送変数とする．
        $\hat{\bm n}_{lg} = \bnabla\phi$，$\kappa_{lg} = \bnabla\cdot\hat{\bm n}_{lg}$ と簡潔に書けるため
        幾何計算精度が高い．しかし移流方程式が非保存形のため，
        数値積分誤差が質量損失として蓄積するという根本的問題がある．

  \item[\textbf{Conservative Level Set（CLS）法}~\cite{OlssonKreiss2005}（\textbf{本稿採用}）]
        $\phi$ を滑らかなヘヴィサイド関数 $\psi = H_\varepsilon(-\phi) \in [0,1]$ に変換し，
        保存形の移流方程式を解く．
        保存形移流による体積管理と距離場に基づく幾何評価を両立させる．
        ただし，幾何量（法線・曲率）の計算に $\phi$ を要するため，
        $\psi \to \phi$ の逆変換はロジット関数 $\phi = \varepsilon\ln((1-\psi)/\psi)$ により
        解析的に求まる（第~\ref{sec:levelset}章参照）．
\end{description}

3手法の特性を表~\ref{tab:method_comparison}に比較する．
CLS 法は質量保存性と幾何精度を同じ枠内で扱える一方，
強変形界面の形状精度は固定格子・閾値化・再初期化の誤差にも律速されるため，
本稿では質量保存軸と形状精度軸を分けて検証する
（§~\ref{sec:zalesak_cls_advection}/§~\ref{sec:single_vortex_reversal}，§~\ref{sec:error_budget}）．

\begin{table}[htbp]
\centering
\caption{界面捕捉 3 手法の特性比較}
\label{tab:method_comparison}
\renewcommand{\arraystretch}{\PaperTableArrayStretch}
\begin{tabular}{lccc}
\hline
特性               & VOF & 標準 LS & CLS（本稿） \\
\hline
質量保存性         & 優  & 劣      & 優 \\
曲率計算精度       & 劣  & 優      & 優 \\
実装の簡潔さ       & 優  & 優      & 可$^*$ \\
界面形状の鮮鋭さ   & 可  & 優      & 可--優 \\
\hline
\end{tabular}

{\footnotesize $^*$ CLS は $\psi \leftrightarrow \phi$ 変換（ロジット関数）・
Ridge--Eikonal 再初期化・保存形フラックス・体積閉鎖の実装が追加で必要なため，
標準 LS に比べやや複雑．逆変換は解析的（Newton 法不要）であり，
本稿の標準再初期化は反復比較経路ではなく 1 ショットの距離再構成として扱う（§~\ref{sec:reinit}参照）．}
\end{table}

% ─────────────────────────────────────────────────────────────────────────────
\subsubsection{寄生流れ抑制手法の先行研究と本稿の位置付け}
% ─────────────────────────────────────────────────────────────────────────────

寄生流れ抑制に対する従来アプローチは大きく二つに分類される．
第一は CSF モデルの枠組み内で演算子整合性を確保する Balanced--Force 条件であり，
Fran{\c{c}}ois ら~\cite{Francois2006} が VOF 枠組みで体系化した．
しかし CSF モデルでは界面遷移層 $\varepsilon$ に起因する近似誤差 $\Ord{\varepsilon^2} \approx \Ord{h^2}$ が
拡散界面モデルの誤差床として残り，演算子整合だけでは除去できない
（§~\ref{sec:balanced_force} に詳述）．

第二は Extension PDE~\cite{Aslam2004} により圧力場を界面越しに滑らかに延長し，
CCD ステンシルが界面近傍でも滑らかな場を参照できるようにするアプローチである．
Desjardins ら~\cite{Desjardins2008} は Ghost-Fluid Method（GFM）と高次 Level Set の組み合わせで
密度不連続を界面で鋭く扱う手法を示しており，界面越し場延長の有効性を裏付けている．

本稿は第一（BF）と第二（界面越し場延長）の両アプローチを統合する立場をとり，
BF の離散整合条件を 7 原則 P-1--P-7（§~\ref{sec:bf_seven_principles}）として系統化した上で，
GFM と Extension PDE の面ジェット統合（HFE，§~\ref{sec:field_extension} で定式化）を CCD 高次精度の枠内で実装する．
具体的には以下の 3 つの柱で寄生流れを抑制する：
\begin{enumerate}[noitemsep]
  \item \textbf{CCD $\Ord{h^6}$ 微分演算子}：
        同一3点ステンシルで $\Ord{h^6}$ の1・2階微分を提供し，
        曲率計算と圧力勾配の両方を高精度化する（第~\ref{sec:CCD}章）．
        ただし $\Ord{h^6}$ は内点ステンシルの設計次数であり，
        境界 BC の打ち切り精度の影響により大域 $L^2$ 漸近では $\Ord{h^{5/2}}$ に低下する点は §~\ref{sec:ccd_def} で詳述する．
  \item \textbf{分相圧力ジャンプ PPE + DC + 面上補正量}：
        高密度比向けの主たる設計戦略として向き付き圧力ジャンプ分相 PPE を採用し，
        各相内で欠陥補正法（DC $k=3$）を適用して高次の圧力補正を得る．
        速度補正・浮力残差・界面移流は同じ面発散・面勾配の組
        $D_fG_A$ と射影後の面速度で閉じ，
        履歴圧力のアフィンジャンプ補正も面加速度として持ち越す．
        一括 CSF + Balanced--Force 一流体経路（§~\ref{sec:balanced_force}）では
        表面張力と圧力勾配に同一の面中心演算子族を使用し，
        分相経路では界面ジャンプ・補正・界面移流の離散位置を一致させることで
        高次の離散化整合性を確保する
        （付録~\ref{sec:ppe_discretization_choice}）．
        検証済み範囲は第~\ref{sec:verification}章でまとめ，
        非一様格子上の移動界面 CLS と長時間・完全連成 NS 移動界面の収束性は
        §~\ref{sec:future_work} の今後の検証対象に残す．
  \item \textbf{Balanced--Force + FCCD 面ジェットによるチェッカーボード抑制}：
        Balanced--Force 演算子整合と FCCD 面ジェット部分系が
        コロケート格子のチェッカーボード不安定をスキーム内部で抑制する
        （§~\ref{sec:checkerboard_solution}；
        運動量移流散逸は UCCD6（§~\ref{sec:uccd6_def}）が担う）．
\end{enumerate}
さらに，高密度比対応の分相 PPE 解法（§~\ref{sec:split_ppe_motivation}；密度比 $\rho_r\le 833$ までの結合閉包系の安定性は §~\ref{sec:density_ratio_sweep} で確認）の基盤技術として，
Hermite 場延長法（HFE，§~\ref{sec:field_extension}）を開発した．
HFE は Extension PDE~\cite{Aslam2004} の定常解を CCD の Hermite データ $(f,f',f'')$ から
仮想時間前進を経由せず直接構成し，界面越し場延長精度を
$\Ord{h}$（風上差分）から大幅に向上させる：1D では $\Ord{h^6}$ を達成し，
2D 斜め界面では完全テンソル版で $\Ord{h^6}$，
混合微分を省略した簡略版では設計次数を保証しない
（§~\ref{sec:U6_split_ppe_dc_hfe} では完全テンソル実装で 2D 低次化が解消する）．
以上により，界面近傍の場延長，圧力補正，速度補正を同じ離散構造で読める定式化を実現する．

なお，寄生流れ低減に対しては Height-Function 曲率推定~\cite{Popinet2009}%
（VOF 法で $\Ord{h^2}$ の曲率精度を達成），CLSVOF 法~\cite{Sussman2000}，
PROST 法~\cite{Renardy2002}%
等の手法も提案されている．
本研究はレベルセット法の枠組み内での高精度化を目指す点でこれらとは異なるアプローチを採る．

% ═══════════════════════════════════════════════════════════════════════════════
\subsection{本稿のアプローチ・構成・読み進め方}
\label{sec:roadmap}
% ═══════════════════════════════════════════════════════════════════════════════

本稿は，単相 Navier--Stokes ソルバーを理解した読者が，
界面追跡・表面張力の離散化・変密度圧力求解を
一貫した演算子の対応関係として読めるように構成する．
導入章では教育的な順序で困難を並べるが，本文の主張は検証済みの研究成果と
未検証範囲の区別に基づいて記述する．

§~\ref{sec:challenges} で述べた4つの困難に対する本稿の手法的対処は，
各困難の説明と対応する「→ 本稿の対処」として同節に整理してある．
ここでは全体像のみ示す：CLS 法による界面追跡，
Ridge--Eikonal/段階 F による幾何プロファイルと体積閉鎖，
FCCD 面ジェット（CLS）と UCCD6（運動量）による移流，
CCD/FCCD/HFE による界面近傍データ構成，
および分相圧力ジャンプ PPE + 欠陥補正法（DC $k=3$）+
射影後の面速度まで同じ面上で共有する補正構造の組み合わせである．

NS 方程式の各項とそれを処理するソルバーの対応関係（演算子マッピング）は，
第~\ref{sec:algorithm}章で 7 ステップ手順として統合し，
各演算子の詳細は第II部の対応章で定式化する．

% ─────────────────────────────────────────────────────────────────────────────
\subsubsection{対象読者と前提知識}
% ─────────────────────────────────────────────────────────────────────────────

本稿を効果的に読み進めるために，以下の前提知識を持っていることが望ましい．
\begin{itemize}
  \item \textbf{流体力学の基礎：}
        非圧縮性 Navier--Stokes 方程式，連続の式，Reynolds 数・Froude 数・Weber 数の概念
  \item \textbf{数値解析の基礎：}
        有限差分法（1次・2次精度の中心差分，前進差分）の概念，
        Taylor 展開，線形代数（連立方程式の反復解法）
  \item \textbf{単相流の数値計算（推奨）：}
        Chorin の Projection 法（圧力--速度分離），スタガード格子または
        コロケート格子の概念，陽的・陰的時間積分の基礎
\end{itemize}
これらに不安がある読者は，まず単相流の標準的な CFD テキストで基礎を固めた上で，
気液二相流の直接数値シミュレーションに関する標準的テキスト
（例：Tryggvason ら~\cite{Tryggvason2011}, Osher--Fedkiw~\cite{OsherFedkiw2003}, Sethian~\cite{Sethian1999}, および VOF/界面追跡の総説~\cite{ScardovelliZaleski1999}）を参照することを推奨する．

% ─────────────────────────────────────────────────────────────────────────────
\subsubsection{各章の概要と依存関係}
% ─────────────────────────────────────────────────────────────────────────────

本稿の章立てと各章の主な内容，および読み進める上での前提章を
表~\ref{tab:chapter_overview}に示す．

\begin{table}[htbp]
\centering
\caption{本稿の章立て・主な内容・前提章の対応}
\label{tab:chapter_overview}
\renewcommand{\arraystretch}{\PaperTableArrayStretch}
\begin{tabular}{clp{68mm}l}
\toprule
章 & タイトル（略称） & 主な内容 & 前提章 \\
\midrule
1  & はじめに
   & 問題背景，失敗例，手法概観，本稿の読み方
   & --- \\
2  & 支配方程式系
   & One-Fluid モデル，CSF 表面張力，無次元化，曲率定義
   & --- \\
3  & CLS 法
   & 保存形界面追跡，再初期化方程式，曲率数値計算
   & 2 \\
4  & CCD 法
   & CCD/DCCD/UCCD6/FCCD と面ジェットの役割分担
   & --- \\
5  & CLS 再初期化
   & Ridge--Eikonal 再初期化，$\phi$ 空間体積閉鎖，CLS の段階的更新手順
   & 2, 3, 4 \\
6  & 方程式項別空間離散
   & CLS $\psi$ 移流，運動量移流，粘性 3 層，物性補間
   & 3, 4, 5 \\
7  & 時間積分
   & CLS は TVD-RK3，NS は EXT2/AB2（IMEX-BDF2）+ IPC，粘性は implicit-BDF2 + DC，時間刻み制約
   & 2, 3, 4, 6 \\
8  & 圧力・速度連成
   & コロケート格子，Helmholtz 分解，Balanced--Force/FCCD 部分系
   & 4, 5 \\
9  & 圧力 Poisson 方程式
   & 分相 PPE 解法，HFE，欠陥補正法，境界条件，精度総括
   & 7, 8 \\
10 & 界面適合格子
   & レベルセット連動格子生成，計量係数，ALE 簡略化，Ridge--Eikonal の非一様格子拡張（D1--D4 数学的定義）
   & 3, 4, 6, 8, 9 \\
11 & ソルバ統合
   & 演算子マッピング，1 タイムステップ手順，純 FCCD DNS 構成（CLS 移流は FCCD 面ジェット，運動量移流は UCCD6）
   & 4--10 \\
12 & コンポーネント検証
   & CCD/FCCD/UCCD6，PPE，非一様格子，Ridge--Eikonal，HFE，BF，時間積分の単体検証
   & 4--10 \\
13 & 統合検証
   & 力学的バランス・保存則・時間空間精度・精度バジェット
   & 12 \\
14 & 多相流物理検証
   & 毛管波・振動液滴・気泡上昇・Rayleigh--Taylor 不安定性の物理検証；長時間・高 We/低 We の定量拡張は §~\ref{sec:future_work} の課題
   & 13 \\
15 & まとめ
   & 技術的貢献の整理，今後の拡張方向
   & --- \\
\bottomrule
\end{tabular}
\end{table}

% ─────────────────────────────────────────────────────────────────────────────
\subsubsection{読み進め方とアルゴリズム全体像}
\label{sec:algo_overview}
% ─────────────────────────────────────────────────────────────────────────────

本稿は，以下の7段階で読み進めると，設計原理から検証範囲までを無理なく追える（表~\ref{tab:chapter_overview}の章構成と対応）．
\begin{enumerate}
  \item \textbf{前提知識の確認：}
        単相 Navier--Stokes 方程式の基礎，Projection 法による圧力--速度分離，有限差分法の基本概念
  \item \textbf{数学的基盤と界面追跡}（第~\ref{sec:governing}--\ref{sec:advection}章）：
        One-Fluid 定式化，CSF モデル（§~\ref{sec:csf}），CLS 法，FCCD（CLS）/UCCD6（運動量）移流スキーム
  \item \textbf{空間離散化スキーム}（第~\ref{sec:CCD}, \ref{sec:advection}, \ref{sec:collocate}, \ref{sec:ppe_solve}, \ref{sec:grid} 章；表~\ref{tab:chapter_overview} の第 4 章・第 6 章・第 8--10 章に相当）：
        CCD 法，CLS/運動量移流，Balanced--Force，HFE，分相 PPE，界面適合格子
  \item \textbf{完全アルゴリズムの統合}：
        第~\ref{sec:algorithm}章で 7 ステップ手順として統合し，各項の離散化は第II部の対応章で確認する
  \item \textbf{コンポーネント検証}（第~\ref{sec:numerical_integrity}章）：
        CCD/FCCD/UCCD6，PPE，非一様格子，界面表現，HFE，BF，時間積分の単体検証
  \item \textbf{統合検証}（第~\ref{sec:verification}章）：
        単体検証で確立した基本構成要素を，単相・二相・非一様・CLS 強変形の統合検証へ接続し，
        合格範囲，条件付き合格，保証範囲外を分離する
  \item \textbf{多相流物理検証}（第~\ref{sec:validation}章）：
        毛管波・振動液滴・気泡上昇・Rayleigh--Taylor 不安定性の物理検証；長時間・高 We/低 We の定量拡張は §~\ref{sec:future_work} の課題
\end{enumerate}

以降の章を読み進める前に，各時間ステップで何が行われるかの実装の地図を把握しておくことが
不可欠である．アルゴリズム全体は（1）界面移流・再初期化・物性更新・曲率/界面ジャンプデータ構成という
界面部分系（ステップ1--4）と，（2）予測・PPE 求解・補正という
流体力学部分系（ステップ5--7）の7ステップで構成される．
両部分系は物性値 $\rho,\mu$ と曲率 $\kappa$ を介して結合する．
各ステップの離散方程式・実装上の注意・収束判定の詳細は，
第~\ref{sec:algorithm}章を入り口として第II部の対応章で述べる．

上記7ステップの関係を図~\ref{fig:algo_flow}に示す．

%%% ── 図：1タイムステップのフロー（高レベル概観版） ──────────────────────
\begin{figure}[htbp]
\centering
\begin{tikzpicture}[
  node distance=5mm and 18mm,
%% 界面部分系用ボックス（青系）
  box/.style={draw=navy!70, fill=lightblue, rounded corners=5pt,
              font=\small\bfseries, inner sep=6pt,
              minimum width=42mm, minimum height=13mm, align=center},
%% 流体力学部分系用ボックス（赤系）
  boxr/.style={draw=red2!70, fill=red2!8, rounded corners=5pt,
               font=\small\bfseries, inner sep=6pt,
               minimum width=42mm, minimum height=13mm, align=center},
  %% 矢印スタイル
  arr/.style={-Stealth, thick, darkgray!80},
  cross/.style={-Stealth, thick, teal2!80, dashed, line width=1.2pt},
  lbl/.style={font=\scriptsize\bfseries, text=darkgray},
  %% 章参照ラベル
  chap/.style={font=\scriptsize\itshape, text=darkgray!65},
]

%% ── 列見出し ──────────────────────────────────────────────────────────
\node[font=\normalsize\bfseries, text=navy]
  (hdL) at (0, 0) {界面部分系};
\node[font=\normalsize\bfseries, text=red2]
  (hdR) at (62mm, 0) {流体力学部分系};

%% ── 左列：界面部分系（ステップ1--4）──────────────────────────────
\node[box, below=8mm of hdL] (s1)
  {1.\ \ CLS 移流\\[1pt]\footnotesize（FCCD 面ジェット）};
\node[box, below=of s1] (s2)
  {2.\ \ 再初期化\\[1pt]\footnotesize（Ridge--Eikonal）};
\node[box, below=of s2] (s3)
  {3.\ \ 物性更新\\[1pt]\footnotesize（$\rho,\,\mu,\,\phi$）};
\node[box, below=of s3] (s4)
  {4.\ \ 曲率・ジャンプデータ\\[1pt]\footnotesize（CCD/HFE）};

%% ── 右列：流体力学部分系（ステップ5--7）──────────────────────────
\node[boxr, right=18mm of s3] (s5)
  {5.\ \ 運動量予測\\[1pt]\footnotesize（Predictor）};
\node[boxr, below=of s5] (s6)
  {6.\ \ 圧力 Poisson 求解\\[1pt]\footnotesize（分相 PPE + DC）};
\node[boxr, below=of s6] (s7)
  {7.\ \ 速度補正\\[1pt]\footnotesize（面上補正）};

%% ── 矢印：左列（界面部分系内）────────────────────────────────────
\draw[arr] (s1) -- (s2);
\draw[arr] (s2) -- (s3);
\draw[arr] (s3) -- (s4);

%% ── 矢印：右列（流体力学部分系内）──────────────────────────────────
\draw[arr] (s5) -- (s6);
\draw[arr] (s6) -- (s7);

%% ── 結合矢印（界面 → 流体力学）──────────────────────────────────────────
\draw[cross] (s3.east)
  -- node[lbl, above, yshift=1pt]{$\rho,\;\mu$} (s5.west);
\draw[cross] (s4.east)
  -- node[lbl, below, yshift=-1pt]{$\kappa$} (s5.south west);
\draw[cross] (s4.east)
  -- node[lbl, above, yshift=1pt]{$J_p,\;a_f$} (s6.west);

%% ── 時間ループ矢印 ───────────────────────────────────────────────────────
\draw[arr]
  (s7.east) -- ++(5mm, 0)
  |- node[lbl, near start, right, xshift=1mm]{\small $t^{n+1} \to t^{n+2}$}
  (s1.east);

\end{tikzpicture}
\caption{1 タイムステップの処理フロー}
\label{fig:algo_flow}
\PaperCaptionNote{概念図．界面部分系（CLS 移流・Ridge--Eikonal 再初期化・物性更新・曲率/ジャンプデータ構成，ステップ1--4）と流体力学部分系（予測・PPE・補正，ステップ5--7）が，物性値 $\rho,\,\mu$，曲率 $\kappa$，および圧力ジャンプ面データを通じて結合している．各ステップの離散方程式・実装詳細は，第~\ref{sec:algorithm}章を入り口として第II部の対応章で述べる．}
\end{figure}
